% =============================================================================
%  Chapter 19 — Explicit Tensor Filters for VTK / ParaView
% =============================================================================

\chapter{Explicit Tensor Filters for VTK / ParaView}
\label{ch:vtk_tensor_filters}

The current VTK contract in \texttt{fall\_n} is intentionally simple:
\begin{itemize}[nosep]
  \item \texttt{*\_mesh.vtu} stores the continuum mesh and nodal fields;
  \item \texttt{*\_gauss.vtu} stores one point per material point and the
        quadrature fields \texttt{qp\_*};
  \item both datasets can carry \texttt{displacement}, and the Gauss cloud
        receives the interpolated material-point displacement explicitly.
\end{itemize}

Within that split workflow, the remaining post-processing problem is not data
transport but \emph{field semantics}.  A raw array named simply
\texttt{stress} or \texttt{strain} is a weak contract: it hides both the
carrier and the tensor convention.  This chapter documents the cleaned-up
naming and the explicit scalar filters now emitted by the exporter.


% ─────────────────────────────────────────────────────────────────────────
\section{Naming Contract}
\label{sec:vtk_tensor_filters_naming}

The exporter distinguishes both \emph{location} and \emph{representation}:
\begin{itemize}[nosep]
  \item \texttt{nodal\_*}: values projected onto mesh nodes;
  \item \texttt{qp\_*}: values evaluated or stored at material points;
  \item \texttt{\_voigt}: raw symmetric-tensor storage in the library's Voigt
        convention;
  \item scalar suffixes such as \texttt{\_xx}, \texttt{\_von\_mises}, or
        \texttt{\_equivalent\_strain}: explicit post-processing views.
\end{itemize}

Examples:
\begin{itemize}[nosep]
  \item \texttt{qp\_stress\_voigt}, \texttt{qp\_stress\_xx},
        \texttt{qp\_stress\_von\_mises};
  \item \texttt{nodal\_strain\_voigt}, \texttt{nodal\_strain\_xx},
        \texttt{nodal\_strain\_equivalent\_strain};
  \item \texttt{qp\_plastic\_strain\_voigt},
        \texttt{nodal\_plastic\_strain\_equivalent\_strain}.
\end{itemize}

This is the main usability improvement.  The user no longer has to infer
whether a field is nodal or quadrature-based, and the exporter no longer asks
ParaView to reinterpret unnamed raw tensors on the fly.


% ─────────────────────────────────────────────────────────────────────────
\section{Stress Filters}
\label{sec:vtk_tensor_filters_stress}

For a symmetric stress tensor $\bsigma$, the exporter emits:
\begin{itemize}[nosep]
  \item component scalars
        $\sigma_{xx}, \sigma_{yy}, \sigma_{zz}, \sigma_{yz}, \sigma_{xz},
        \sigma_{xy}$;
  \item trace
        \[
          \tr(\bsigma) = \sigma_{xx} + \sigma_{yy} + \sigma_{zz};
        \]
  \item mean / hydrostatic stress
        \[
          \sigma_m = \sigma_h = \frac{1}{d}\tr(\bsigma);
        \]
  \item pressure
        \[
          p = -\sigma_m;
        \]
  \item deviatoric norm
        \[
          \|\dev \bsigma\| = \sqrt{\dev \bsigma : \dev \bsigma};
        \]
  \item second deviatoric invariant
        \[
          J_2 = \frac{1}{2}\,\dev \bsigma : \dev \bsigma;
        \]
  \item von Mises / Beltrami--Haigh equivalent stress
        \[
          \sigma_{\mathrm{vm}}
          =
          \sigma_{\mathrm{BH}}
          =
          \sqrt{\frac{3}{2}\,\dev \bsigma : \dev \bsigma};
        \]
  \item triaxiality
        \[
          \eta = \frac{\sigma_m}{\sigma_{\mathrm{vm}}};
        \]
  \item octahedral shear stress
        \[
          \tau_{\mathrm{oct}} = \sqrt{\frac{2}{3}J_2}.
        \]
\end{itemize}

The exporter writes both \texttt{von\_mises} and
\texttt{beltrami\_haigh}.  In this context they are numerically the same
$J_2$-based equivalent stress, but exposing both names removes friction for
users who search for one label or the other in ParaView.


% ─────────────────────────────────────────────────────────────────────────
\section{Strain Filters}
\label{sec:vtk_tensor_filters_strain}

The strain arrays are stored in engineering Voigt notation:
\[
  \{\varepsilon_{xx}, \varepsilon_{yy}, \varepsilon_{zz},
    \gamma_{yz}, \gamma_{xz}, \gamma_{xy}\},
  \qquad
  \gamma_{ij} = 2\varepsilon_{ij}.
\]

The exporter therefore converts the off-diagonal entries back to tensor shear
components before evaluating invariants.  The derived strain filters are:
\begin{itemize}[nosep]
  \item component scalars in the stored engineering convention;
  \item trace / volumetric strain
        \[
          \varepsilon_v = \tr(\bvarepsilon);
        \]
  \item deviatoric norm
        \[
          \|\dev \bvarepsilon\|
          =
          \sqrt{\dev \bvarepsilon : \dev \bvarepsilon};
        \]
  \item equivalent strain
        \[
          \varepsilon_{\mathrm{eq}}
          =
          \sqrt{\frac{2}{3}\,\dev \bvarepsilon : \dev \bvarepsilon}.
        \]
\end{itemize}

The same reconstruction is applied to plastic-strain-like quantities, which is
why arrays such as \texttt{qp\_plastic\_strain\_xx} and
\texttt{nodal\_plastic\_strain\_equivalent\_strain} are available without any
ParaView-side calculator.


% ─────────────────────────────────────────────────────────────────────────
\section{Why the Exporter Owns These Scalars}
\label{sec:vtk_tensor_filters_why}

There are two architectural reasons to compute the scalar filters inside
\texttt{fall\_n}:
\begin{enumerate}[nosep]
  \item the library owns the tensor convention and therefore should own the
        conversion rules;
  \item explicit scalar arrays are stable extension points for downstream
        scripts, regression tests, dashboards, and future post-processing
        modules.
\end{enumerate}

This is compatible with the HPC direction of the library.  The additional work
is restricted to post-processing, while the FE hot path remains unchanged.


% ─────────────────────────────────────────────────────────────────────────
\section{Recommended ParaView Usage}
\label{sec:vtk_tensor_filters_paraview}

For smooth contours on the continuum mesh:
\begin{enumerate}[nosep]
  \item load \texttt{*\_mesh.vtu};
  \item if desired, apply \texttt{Warp By Vector} with \texttt{displacement};
  \item colour by \texttt{nodal\_stress\_von\_mises},
        \texttt{nodal\_strain\_xx}, or any other explicit nodal scalar.
\end{enumerate}

For material-point inspection:
\begin{enumerate}[nosep]
  \item load \texttt{*\_gauss.vtu};
  \item apply \texttt{Warp By Vector} directly using \texttt{displacement};
  \item colour by \texttt{qp\_stress\_von\_mises},
        \texttt{qp\_stress\_xx}, \texttt{qp\_strain\_equivalent\_strain}, or
        similar explicit scalars;
  \item use \texttt{Points} or \texttt{Point Gaussian} representation as
        appropriate.
\end{enumerate}

The raw \texttt{\_voigt} arrays remain in the files for traceability and
advanced tooling, but they are no longer the primary user-facing
representation.

